%! AP Statistics — Unit 4, Lesson 4.11 — Warm-Up KEY
\documentclass[10pt]{article}
\usepackage{apstats-article}
\usepackage{apstats-key}

\begin{document}

\pageheader{Unit 4, Lesson 4.11}{Warm-Up --- Answer Key}
\namedateperiod

\vspace{0.1in}

% ── SECTION 1: Spiral Review ──────────────────────────────────────────────────
{\large\bfseries\color{navy} Part 1 \quad Spiral Review}

\vspace{0.08in}

\begin{enumerate}

    \item Let $X \sim B(10, 0.3)$.
    \begin{enumerate}[label=(\alph*), itemsep=4pt]
        \item Verify that this is a binomial setting by listing all four BINS conditions.\\
        \ansline{B: each trial is a success (prob.\ 0.3) or a failure.}\\[1pt]\ansline{I: trials are independent of one another.}\\[1pt]\ansline{N: fixed number of trials, $n = 10$.}\\[1pt]\ansline{S: same $p = 0.3$ on every trial.}
        \item Compute $P(X = 4)$. Show the formula setup.\\
        \ansline{$P(X=4) = \binom{10}{4}(0.3)^4(0.7)^6$}\\[1pt]\ansline{$= 210(0.0081)(0.1176) \approx 0.2001$}
    \end{enumerate}

    \vspace{0.1in}

    \item Recall the mean of a \emph{general} discrete random variable:
    $\mu_X = \sum x_i \cdot P(X = x_i)$.

    The table below shows part of the distribution of $X \sim B(3, 0.5)$.
    Complete the table and find $\mu_X$ using the general formula.

    \vspace{4pt}
    \renewcommand{\arraystretch}{1.5}
    \begin{tabularx}{\linewidth}{@{} >{\bfseries}p{1.4cm} X X X X @{}}
    $x$ & 0 & 1 & 2 & 3 \\
    \hline
    $P(X=x)$ & \ans{0.125} & \ans{0.375} & \ans{0.375} & \ans{0.125} \\
    \hline
    $x \cdot P(X=x)$ & \ans{0} & \ans{0.375} & \ans{0.750} & \ans{0.375} \\
    \end{tabularx}

    \vspace{4pt}
    $\mu_X = \sum x_i P(x_i) = $ \ans{0 + 0.375 + 0.750 + 0.375 = 1.5}

    \vspace{0.1in}

    \item Recall that for any discrete random variable, $\sigma_X = \sqrt{\sum(x_i - \mu_X)^2 P(x_i)}$.
    Based on your answer to Problem 2, what would you \emph{predict} for $\mu_X$ using
    the formula $np$? Do your two answers agree?\\
    \ansline{$np = 3(0.5) = 1.5$.}\\[1pt]\ansline{Yes --- both methods give $\mu_X = 1.5$; $np$ is the shortcut.}

\end{enumerate}

\vspace{0.08in}

% ── SECTION 2: Looking Ahead ──────────────────────────────────────────────────
{\large\bfseries\color{navy} Part 2 \quad Looking Ahead}

\vspace{0.08in}

\noindent Today we replace the long summation in Problem 2 with two simple formulas:
$\mu_X = np$ and $\sigma_X = \sqrt{np(1-p)}$.

\vspace{0.1in}

\begin{tcolorbox}[colback=greenbg, colframe=greenacc, boxrule=0.8pt, arc=2mm,
    left=4mm, right=4mm, top=3mm, bottom=3mm]
    \small
    \begin{itemize}[leftmargin=1.3em, itemsep=4pt]
  \item $\mu_X = np$ is the expected (average) number of successes per experiment.
  \item Always interpret $\mu_X$ with a sentence that references the context.
  \item \texttt{binomcdf} gives cumulative probabilities --- useful when the event
        involves a range of values like ``at most 5'' or ``at least 3.''
    \end{itemize}
\end{tcolorbox}

\vspace{0.1in}

\noindent\textcolor{slate}{\small Write one question you have about today's lesson:}\\[4pt]
\ansline{Student responses will vary.}\\[2pt]\ansline{}

\end{document}
